
\documentclass[11pt,a4paper]{article}

%% PACHETE MATEMATICE
\usepackage{amsmath,amssymb,amsfonts}
\usepackage{amsthm}
\usepackage{mathpazo}
\usepackage{eulervm}
\usepackage{enumerate}
\usepackage{color}
\usepackage{graphicx}
\usepackage[all]{xy}
\usepackage{xcolor}
\usepackage{framed}
\usepackage[unicode]{hyperref}
\setcounter{MaxMatrixCols}{20}
\usepackage{multicol}
%\usepackage[romanian]{babel}


%% ASPECT
\linespread{1.05}
\usepackage[T1]{fontenc}
\usepackage[utf8x]{inputenc}
\usepackage[margin=2cm]{geometry}
% \usepackage[color]{showkeys}
\usepackage{anyfontsize}
\usepackage{titlesec}
\titleformat*{\section}{\large\bfseries}

\usepackage{fancyhdr}
\pagestyle{fancy}
\rhead{\small \color{gray} Notițe de Adrian Manea}
\lhead{\small \color{gray} M1-ACS, 2017---2018, M. Olteanu}


\usepackage[autostyle = false, style = german]{csquotes}
\usepackage[romanian]{babel}
\newcommand\qq{\enquote}

%% COMENZILE MELE
\renewcommand*{\proofname}{Dem.:}
\newcommand\bb{\mathbb}
\newcommand\dr{\mathrm}
\newcommand\kal{\mathcal}
\newcommand\fk{\mathfrak}
\newcommand\op{\oplus}
\newcommand\ot{\otimes}
\newcommand\ds{\displaystyle}
\newcommand\id{\indent}
\newcommand\nid{\noindent}
\newcommand\seq{\subseteq}
\newcommand\sneq{\subsetneq}
\newcommand\speq{\supseteq}
\newcommand\spneq{\supsetneq}
\newcommand\rar{\rightarrow}
\newcommand\Rar{\Rightarrow}
\newcommand\xrar{\xrightarrow}
\newcommand\lar{\leftarrow}
\newcommand\Lar{\Leftarrow}
\newcommand\xlar{\xleftarrow}
\newcommand\lrar{\leftrightarrow}
\newcommand\Lrar{\Leftrightarrow}
\newcommand\ideal{\trianglelefteq}
\newcommand\ai{\text{ a.\^{i}. }}
\newcommand\sk{(\textit{Schi\c{t}\u{a}}) }
\newcommand\rhup{\rightharpoonup}
\newcommand\rhdn{\rightharpoondown}
\newcommand\lhup{\leftharpoonup}
\newcommand\lhdn{\leftharpoondown}
\newcommand\vid{\emptyset}
\newcommand\wh{\widehat}
\newcommand\ol{\overline}
\newcommand\NN{\mathbb{N}}
\newcommand\RR{\mathbb{R}}
\newcommand\ZZ{\mathbb{Z}}
\newcommand\QQ{\mathbb{Q}}
\newcommand\CC{\mathbb{C}}

%% STILURILE MELE
\newtheoremstyle{dotless}{}{}{\itshape}{}{\bfseries}{:}{ }{}

\theoremstyle{dotless}
\newtheorem{theorem}{Teorem\u{a}}[section]
\newtheorem{proposition}{Propozi\c{t}ie}[section]
\newtheorem{lemma}{Lem\u{a}}[section]
\newtheorem{corollary}{Corolar}[section]
\newtheorem{exercise}{Exerci\c{t}iu}

\newtheoremstyle{dotlessdr}{}{}{}{}{\bfseries}{:}{ }{}
\theoremstyle{dotlessdr}
\newtheorem{example}{Exemplu}[section]
\newtheorem{definition}{Defini\c{t}ie}[section]
\newtheorem{remark}{Observa\c{t}ie}[section]




\begin{document}

\begin{center}
  {\large\textbf{Seminar 13b \\ Exerciții recapitulative}}
\end{center}

\vspace{1cm}

\textbf{MODEL 1}

1. Să se determine maximul și minimul funcției $f(x, y) = x^2 + y^2 - 3x - 2y + 1$
pe mulțimea $K = \{(x, y) \in \RR^2 \mid x^2 + y^2 \leq 1 \}$.

\vspace{1cm}

2. Să se calculeze $\ds\int_0^\infty \frac{x^{\frac{1}{4}}}{(x + 1)^2} dx$.

\vspace{1cm}

3. Să se calculeze integrala:
\[
  \iint_D 1 + \sqrt{x^2 + y^2} dx dy,
\]
unde $D = \{(x, y) \in \RR^2 \mid x^2 + y^2 - 3 \leq 0, x \geq 0 \}$.

\vspace{1cm}

4. Să se calculeze aria paraboloidului $z = x^2 + y^2, z \in [0, h]$.

\vspace{1cm}

5. Să se calculeze integrala de suprafață:
\[
  \int_\Sigma x^2 dy \wedge dz - 2xy dz \wedge dx + z^3 dx \wedge dy,
\]
unde $\Sigma : x^2 + y^2 + z^2 = 9$.

\vspace{2cm}

\textbf{MODEL 2}

1. Să se determine punctele de extrem local ale funcției:
\[
  f(x, y) = 3xy^2 - x^3 - 15x - 36y + 9.
\]

\vspace{1cm}

2. Să se studieze convergența integralei
$\ds\int_1^\infty \frac{\ln x}{\sqrt{x^3 - 1}} dx$.

\vspace{1cm}

3. Să se calculeze fluxul cîmpului de vectori $\vec{r} = x\vec{i} + y\vec{j} + z \vec{k}$
prin suprafața de ecuație $z = x^2 + y^2, z \in [0, 5]$.

\vspace{1cm}

4. Fie $P, Q, R : \Omega = \{(x, y, z) \in \RR^3 \mid y > 0, z \geq 0 \} \to \RR$,
cu:
\begin{align*}
  P &= x^2 - yz + \frac{y}{x^2 + y^2} \\
  Q &= y^2 - zx - \frac{x}{x^2 + y^2} \\
  R &= z^2 - xy.
\end{align*}
Notăm $\omega = Pdx + Qdy + Rdz$. Să se calculeze $\ds\int_\Gamma \omega$, unde
$\Gamma$ este un drum parametrizat ce unește punctele $A(1, 1, 0)$ și $B(-1, 1, 0)$.

\vspace{1cm}

5. Să se calculeze, folosind formula Green-Riemann, integrala
$\ds\int_\Gamma xydx + \frac{x^2}{2} dy$, pe mulțimea $\Gamma = \Gamma_1 \cup \Gamma_2$,
unde:
\begin{align*}
  \Gamma_1 &= \{(x, y) \mid x^2 + y^2 = 1, x \leq 0 \leq y \} \\
  \Gamma_2 &= \{(x, y) \mid x + y = -1, x \leq 0, y \leq 0 \}.
\end{align*}

\vspace{2cm}

\textbf{MODEL 3}

1. Să se calculeze punctele de extrem și valorile extreme pentru funcția:
\[
  f(x, y) = xy(x + y - 2),
\]
pe mulțimea $K = \{(x, y) \in \RR^2 \mid x + y \leq 3, y \geq 0\}$.

\vspace{1cm}

2. Să se calculeze integrala:
\[
  I = \int_0^1 \ln^p \frac{1}{x} dx, p > -1.
\]

\vspace{1cm}

5. Să se calculeze, folosind formula lui Stokes, integrala curbilinie
$\int_\Gamma \alpha$, pentru $\alpha = (y - z) dx + (z - x)dy + (x - y) dz$,
iar mulțimea $\Gamma: z = x^2 + y^2, z = 1$.

\vspace{1cm}

4. Fie $\vec{V} = (x^2 + y - 4) \vec{i} + 3xy\vec{j} + (2xz + z^2) \vec{k}$ și
emisfera $\Sigma: x^2 + y^2 + z^2 = 16, z \geq 0$.

Să se calculeze fluxul cîmpului $\nabla \times \vec{V} = \mathrm{rot} \vec{V}$ prin
$\Sigma$, orientată cu normala exterioară la sferă.

\vspace{1cm}

5. Fie $\alpha = \frac{x - y}{x^2 + y^2} dx + \frac{x + y}{x^2 + y^2} dy$.
Să se calculeze integrala curbilinie $\int_\Gamma \alpha$, unde $\Gamma$ este triunghiul
cu vîrfurile $A(0, 3), B(0, 0), C(4, 0)$.

\end{document}